\par\noindent\textbf{(1)} 总激发数算符 \(N_{\rm exc}=a^\dagger a+\sigma_+\sigma_-\)，由 \([a^\dagger a,a]=-a\)、\([\sigma_+\sigma_-,\sigma_+]=\sigma_+\) 得 \([N_{\rm exc},a\sigma_+]=0\)，共轭项亦为零，所以 \([H,N_{\rm exc}]=0\)。激发数为 \(n+1\) 的二维不变子空间由 \(|e,n\rangle,|g,n+1\rangle\) 张成。
\par\noindent\textbf{(2)} 在该基底减去中心能量 \(\hbar\omega_c(n+1/2)\)，剩余 \(2\times2\) 矩阵为 \(\frac\hbar2\begin{pmatrix}\Delta&2g\sqrt{n+1}\\2g\sqrt{n+1}&-\Delta\end{pmatrix}\)。两个极化能为 \(E_{n,\pm}=\hbar\omega_c(n+1/2)\pm\frac\hbar2\sqrt{\Delta^2+4g^2(n+1)}\)，混合角满足 \(\tan(2\theta_n)=2g\sqrt{n+1}/\Delta\)。
\par\noindent\textbf{(3)} 初态为 \(|e,n\rangle\) 时，两个混合本征态之间能级差为 \(\hbar\Omega_n\)，其中 \(\Omega_n=\sqrt{\Delta^2+4g^2(n+1)}\)。精确激发态占有率 \(P_e(t)=1-\frac{4g^2(n+1)}{\Omega_n^2}\sin^2(\Omega_nt/2)\)。共振 \(\Delta=0\) 时 \(P_e(t)=\cos^2(g\sqrt{n+1}\,t)\)，可完全转移激发到 \(|g,n+1\rangle\)。
\par\noindent\textbf{(4)} 在 \(|\Delta|\gg2g\sqrt{n+1}\) 的色散区，展开 \(\sqrt{\Delta^2+4g^2(n+1)}=|\Delta|+2g^2(n+1)/|\Delta|+O[g^4(n+1)^2/|\Delta|^3]\)。取带符号 \(\chi=g^2/\Delta\)，有效Hamilton量包含 \(\hbar\chi(a^\dagger a+1/2)\sigma_z\)，给原子态依赖的腔频移 \(\pm\chi\)；临界光子数尺度 \(n_{\rm crit}\sim\Delta^2/(4g^2)\)。
